% !TEX root = misc-gadget.tex
% title: 杂记：gadget 向量

\subsection{【补丁｜2026-976】gadget 向量与安全界的修订稿}\label{sec:gadget-vector}

\begin{patchnote}{2026-976 与笔记勘误｜本篇正文整体替换稿，包含必要的原有基础公式}
【初版笔记】统一使用 $L=\kappa+2\lambda_s$，要求 gadget 提前固定，
并把 LHL 界误用于单个选择位的猜测成功概率。
标作「2026-976 相关笔记勘误」的段落是阅读该文后对原笔记的纠错，并非该文提出的新协议步骤。
本稿分别标出论文更新与原笔记中的解释性勘误，不能把本篇全部基础公式都理解为论文新增。

【补丁写法】
gadget 向量把 OT 的选择比特串压成一个标量，同时把逐行乘积份额聚合成标量域上的 MtA。
本文记标量模数为 $n$，位宽为 $\kappa=\lceil\log_2n\rceil$，OT 数量为 $L$。
计算与统计安全参数分别为 $\lambda_c,\lambda_s$，与模数位宽分开使用。

DKLS23 用随机权重 $\boldsymbol g\in\mathbb Z_n^L$ 替代朴素的二进制权重。
旧参数 $L=\kappa+2\lambda_s$ 足以描述独立均匀选择串的基本 LHL 估计，
但不能直接覆盖恶意 Sender 在看到 gadget 后施加的选择约束。
本篇按 \href{https://eprint.iacr.org/2026/976.pdf}{Asharov，2026/976，第 4～5 节}
更新选择时机、参数与适用范围；完整协议见 \hyperref[sec:rvole:parameters-model]{RVOLE 完全版}。

\begin{equation}\label{eq:misc-gadget:1}
2^j\;\longmapsto\;g_j,\qquad
L\text{ 按 gadget 的生成时机与安全参数选取。}
\end{equation}

形式上，逐行 OT 计算完成后，双方分别执行
\begin{equation}\label{eq:misc-gadget:2}
\hat x_b=\sum_{j=0}^{L-1}g_j\beta_j,\qquad
z_a=-\sum_{j=0}^{L-1}g_j\alpha_j^0,\qquad
z_b=\sum_{j=0}^{L-1}g_jt_j.
\end{equation}
若逐行关系为 $t_j=\alpha_j^0+\beta_jw$，则 $z_a+z_b=w\hat x_b$。
这条正确性等式与 gadget 何时生成无关，只要求双方使用同一向量。
恶意安全还取决于逐行一致性检查与随机性条件，不能仅凭该代数等式推出。

\subsubsection{【补丁｜2026-976 相关笔记勘误】选择串与聚合标量}

Receiver 的选择串是
\begin{equation}\label{eq:misc-gadget:4}
\boldsymbol\beta=(\beta_0,\ldots,\beta_{L-1})\in\{0,1\}^L.
\end{equation}
它有 $L$ 比特，而 $\hat x_b=\langle\boldsymbol g,\boldsymbol\beta\rangle$ 是 $\mathbb Z_n$ 中的一个标量。
多个不同的选择串可以得到同一个标量，这给逐行一致性检查保留了隐藏选择位的余地。
但「有很多原像」不等于「每个选择位都不可预测」，原像的数量也不直接给出恶意安全界。

朴素二进制权重在取模前就是整数的二进制编码。
取模后，当 $n<2^\kappa$ 时该映射不是严格双射，因此也不能把模 $n$ 的版本称为双射。
随机 gadget 通过增加选择串长度并随机化权重，提供更充足的熵提取空间。

\subsubsection{【补丁｜2026-976 相关笔记勘误】普通 leftover hash lemma 保证什么}

设 $\boldsymbol B$ 来自一个与 $\boldsymbol G$ 独立的源，最小熵至少为 $h$，
而 $\boldsymbol G$ 均匀取自 $\mathbb Z_n^L$。
对不同的二进制串，其内积相等的概率为 $1/n$，所以这些内积映射构成通用哈希族。
LHL 给出
\begin{equation}\label{eq:misc-gadget:5}
\operatorname{SD}\bigl(
  (\boldsymbol G,\langle\boldsymbol G,\boldsymbol B\rangle),
  (\boldsymbol G,U_{\mathbb Z_n})\bigr)
\le \frac12\sqrt{\frac n{2^h}}.
\end{equation}
若 $\boldsymbol B$ 在全部 $L$ 比特串上均匀分布，则 $h=L$。
此时取 $L=\kappa+2\lambda_s$，右侧不大于 $2^{-\lambda_s-1}$。
这是原始独立源的估计，尚未处理恶意 Sender 的选择性失败。

该式比较的是「权重与聚合输出」的联合分布，并不声称输出与原选择串独立。
例如，输出确定为原选择串的函数，本来就携带关于它的约束。
同样，猜一个均匀比特的成功概率是 $1/2$；不能把 $2^{-\lambda_s}$ 写成猜该比特的成功概率。

\subsubsection{【补丁｜2026-976】提前公开 gadget 后，为什么旧参数不足}

在 RVOLE 中，恶意 Sender 可以让一些 OT 行使用不一致的等效输入。
一致性检查可能仅在若干 Receiver 选择位恰好符合它的猜测时通过，这就是选择性失败。
条件化到「检查通过」后，选择串落入受约束的源，分布已不再是原来的独立均匀源。

若 $\boldsymbol g$ 在协议开始时就公开，Sender 可以看着它选择要约束哪些位以及这些位的取值。
于是，条件化后的源可以依赖 $\boldsymbol g$。
普通 LHL 对每个预先固定的高熵源给出保证，不能直接交换量词，
将其改写成「这个已公开的 $\boldsymbol g$ 对所有事后选择的源都足够好」。

2026/976 第 5 节分别分析三种游戏，区分源与 gadget 独立、有限次尝试 gadget，
以及公开 gadget 后自适应固定部分选择位。
第三种情况下，需要对协议允许的选择约束族统一控制坏事件。
这导致更大的 OT 数量充分界；论文给出的是保守界，不是已经证明最优的最小参数。

\subsubsection{【补丁｜2026-976】三种时序及其实施}

【Variant I】Receiver 在校验 Sender 消息后，才均匀选择 $\boldsymbol g$ 并发送。
这样，Sender 形成选择约束时还不知道 gadget，可以使用独立源的分析。
代价是增加一条 Receiver 消息，标量 $\hat x_b$ 也到这个阶段才产生；充分条件为
$L\ge\kappa+3\lambda_s$。

【Variant II】双方从 Sender 的最后一条消息派生 gadget。
沿用 RVOLE 的记号，令最后消息的摘要 $M_2=\mathrm{RO}(\texttt{"rvole-msg"},\tilde a,\eta,\mu)$，再计算
\begin{equation}
\boldsymbol g:=\mathrm{RO}(\mathrm{sid},\texttt{"gadget-II"},T_{\mathrm{OT}},M_2)
  \in \mathbb Z_n^L.
\end{equation}
Sender 可以尝试多个候选消息并查询随机预言机，参数须计入这种研磨能力。
论文给出的充分条件为 $L\ge\kappa+2\lambda_c+3\lambda_s$。
该方案省去单独发送 gadget 的轮次，但 Receiver 仍要收到最后消息后才能得到 $\hat x_b$。

【Variant III】双方在协议开始时生成
\begin{equation}\label{eq:misc-gadget:gvec}
\boldsymbol g:=\mathrm{RO}(\mathrm{sid},\texttt{"gadget-III"})
  \in \mathbb Z_n^L.
\end{equation}
这保留了多方 DKLS23 所需的早期 Receiver 标量，充分条件为
\begin{equation}
L\ge\kappa+7\lambda_s+2\log_2(\lambda_s+1)
       +2\lambda_s\log_2(eL/\lambda_s).
\end{equation}
这里 $e$ 为自然常数，所有对数以 2 为底。
随机预言机到域元素的实现采用 XOF 与拒绝采样，不能直接把短哈希取模称为均匀采样。
各有向实例使用唯一会话号，参数、协议版本和身份都纳入无歧义编码。

\subsubsection{【补丁｜2026-976】与 RVOLE 和签名编排的关系}

前两种方案都把聚合标量的产生推迟到 Receiver 收到 Sender 消息之后。
多方签名第二轮要发送 $\psi_{i,j}=\phi_i-\beta_{j,i}$，因此不能直接换成较晚产出 $\beta_{j,i}$ 的方案。
本笔记的多方流程采用 III；两方每个方向只有一个对端，可采用 II 并省去这层盲化值对齐。

此外，RVOLE 一致性挑战必须约束对应 OT 阶段。
若 Endemic OT 允许恶意 Sender 自选输入，仅由修正矩阵生成挑战可能允许事后调整 OT 输入来抵消检查。
因此，本篇引用的参数与协议结论还要求
\hyperref[sec:rvole:parameters-model]{Sender Random OT 或具体 transcript 输入绑定条件}，
不能把 gadget 参数修订当作这个独立问题的修补。

\subsubsection{【补丁｜2026-976】代价与参数核对}

以 secp256k1、$\kappa=256$、$\lambda_c=\lambda_s=128$ 为例，
I、II、III 的上述充分用量分别为 640、896、2656 个 OT。
为了配合现有 SoftSpoken 的 128 列分块，默认 III 取 $L=2688$，
并将扩展列数设为 $L'=2816$，真实块数设为 $M=21$。

RVOLE 的一个校验负载还要求 $\log_2 n\ge\lambda_c+2\log_2L$，上述参数满足该条件。
扩展数量增加不改变 PPRF 树深 4、树数 64 或 Base OT 数 256，
而是延长叶子伪随机流及后续矩阵的列数。
这些参数分别服务于其协议组件，不能据此独立断言整套 ECDSA 达到了同一统计安全等级。
\end{patchnote}
